\section{Performance Optimizations}


\subsection{Flash Attention}

Flash Attention~\citep{dao2023flashattention} reduces attention memory from $O(N^2)$ to $O(N)$ via tiling:

\begin{verbatim}
# Standard attention: materialize full NxN matrix
Q @ K^T @ V  # Memory: O(N^2)

# Flash Attention: block-sparse computation
for block_i in range(num_blocks):
    for block_j in range(num_blocks):
        compute_attention_block(Q[i], K[j], V[j])
        # Memory: O(block_size^2)
\end{verbatim}

Flash Attention 2 adds:
\begin{itemize}
\item Reduced shared memory usage
\item Better parallelism across SMs
\item 2-3x speedup over Flash Attention 1
\end{itemize}

Gym enables Flash Attention automatically when:
\begin{itemize}
\item GPU supports BF16 or FP16
\item Sequence length > 512 tokens
\item \texttt{transformers >= 4.35.0}
\end{itemize}

\subsection{Liger Kernel}

Liger~\citep{hsu2024liger} optimizes common LLM operations:

\begin{itemize}
\item \textbf{Fused Cross-Entropy}: Combines softmax + NLL loss in single kernel
\item \textbf{Fused RMSNorm}: Layer norm without mean subtraction
\item \textbf{Fused RoPE}: Rotary position embeddings
\item \textbf{Fused SwiGLU}: Gated activation functions
\end{itemize}

Memory savings:
\begin{verbatim}
Standard:   [Activation] -> [Softmax] -> [CE Loss]
            Memory: 3x hidden_dim

Liger:      [Fused Op]
            Memory: 1x hidden_dim (67% reduction)
\end{verbatim}

\subsection{Unsloth}

Unsloth~\citep{unsloth2024} achieves 2x speedup via:
\begin{itemize}
\item Manual attention kernel implementation
\item LoRA-specific fused operators
\item Automatic mixed precision (AMP) optimization
\item Gradient accumulation without extra memory
\end{itemize}

Benchmarks (7B model, A100 40GB):
\begin{table}[h]
\centering
\begin{tabular}{lcc}
\toprule
Configuration & Tokens/sec & Memory (GB) \\
\midrule
Baseline & 1450 & 38.2 \\
Flash Attention 2 & 2100 & 36.5 \\
+ Liger & 2400 & 34.1 \\
+ Unsloth & 2900 & 32.8 \\
\bottomrule
\end{tabular}
\caption{Cumulative speedup from optimization stack.}
\end{table}

\subsection{Quantization}

\subsubsection{Post-Training Quantization (PTQ)}

Gym supports:
\begin{itemize}
\item \textbf{GPTQ}~\citep{frantar2023gptq}: Per-column quantization with Hessian approximation
\item \textbf{AWQ}~\citep{lin2023awq}: Activation-aware weight quantization preserving salient channels
\item \textbf{SmoothQuant}~\citep{xiao2023smoothquant}: Migrates difficulty from activations to weights
\end{itemize}

\subsubsection{Quantization-Aware Training (QAT)}

QLoRA performs QAT implicitly:
\begin{itemize}
\item Base model quantized to 4-bit NF4
\item LoRA adapters trained in FP16/BF16
\item Gradients backpropagate through quantization
\end{itemize}

\textbf{NormalFloat4 (NF4)}: Information-theoretically optimal 4-bit format:
\begin{equation}
\text{NF4} = \{\text{quantiles of } \mathcal{N}(0,1) \text{ at intervals of } 1/16\}
\end{equation}

\subsection{Mixed Precision Training}

Gym automatically enables AMP (Automatic Mixed Precision):
\begin{itemize}
\item Forward pass: FP16/BF16
\item Gradient accumulation: FP32
\item Optimizer states: FP32
\item Loss scaling: Dynamic (FP16) or none (BF16)
\end{itemize}

BF16 preferred over FP16 when available (Ampere+ GPUs):
\begin{itemize}
\item Same dynamic range as FP32 (no loss scaling needed)
\item Fewer NaN/Inf issues
\item Slightly lower precision acceptable for LLMs
\end{itemize}

